% =========================
%  Insulation Coordination
% =========================
\intermezzo{Withstand Voltage}{}


\begin{frame}{Withstand voltage definition process}
  \begin{columns}[c,onlytextwidth]
    \begin{column}{0.5\textwidth}
      \begin{enumerate}
        \setlength{\itemsep}{0.4em}
        \item Determination of the \\representative overvoltages $U_{rp}$
        \item Determination of the co-ordination \\withstand voltages $U_{cw}$. \\For HVDC it can be assumed:\\
              \quad$U_{cw} = U_{rp}$
        \item Determination of the required \\ withstand voltages $U_{rw}$
        \item Determination of the specified \\ withstand voltages $U_w$
      \end{enumerate}
    \end{column}
    \begin{column}{0.5\textwidth}
      \centering
      \topoimg[width=\linewidth,height=0.8\textheight,keepaspectratio]{\icfig{m08_image17}}\\[0.3em]
      {\color{atGray}\footnotesize IEC 60071-11:2022.}
    \end{column}
  \end{columns}
  \vspace{0.2em}
\end{frame}

\begin{frame}{Representative overvoltages ($U_{rp}$)}
  \begin{block}{IEC 60071-11:2022}
    ``... representative overvoltages are equal to protection levels of
    arresters for directly protected equipment.''
  \end{block}
  \begin{itemize}
    \setlength{\itemsep}{0.4em}
    \item The statement is misleading: the arrester protection level is
          defined by the arrester design and the protection-coordination
          current amplitude and shape (lightning, switching).
    \item Representative overvoltages can be determined by assuming maximum
          system currents and switching impulse, or by transient-simulation
          analysis of relevant faults.
    \item Calculation and simulation results without any margin are
          considered as the representative overvoltage $U_{rp}$.
  \end{itemize}
\end{frame}

\begin{frame}{Required withstand voltages ($U_{rw}$)}
  
  $U_{rw}$ is defined as:\\
  {\centering $U_{rw}=K_a \cdot K_s \cdot U_{rp}$\par}
  \vspace{1em}
  
  $K_a$ Ambient condition correction factor
  \begin{itemize}
    \item Altitude correction
    \item Temperature / humidity correction for indoor installations
  \end{itemize}
  \vspace{1em}
  
  $K_s$ Safety factor
  \begin{itemize}
    \item Altitude correction up to 1000m
    \item allowance for ageing of insulation
    \item allowance for changes in arrester characteristics
    \item allowance for dispresion in the product quality
  \end{itemize}

\end{frame}


\begin{frame}{Required withstand voltages ($U_{rw}$)}
      Definition safety factor $K_s$\\
      \centering      
      \topoimg[width=\linewidth,height=0.66\textheight,keepaspectratio]{\icfig{m08_image19}}\\
      {\color{atGray}\footnotesize IEC 60071-11:2022}
\end{frame}

\begin{frame}{Specified withstand voltages ($U_w$)}
  \begin{columns}[c,onlytextwidth]
    \begin{column}{0.54\textwidth}
      \begin{itemize}
        \setlength{\itemsep}{0.4em}
        \item Values are equal to or higher than $U_{rw}$
        \item For AC equipment they correspond to the standard values in
              IEC 60071-1
        \item For HVDC equipment they are rounded up to convenient values\\
              (in most cases IEC 60071-1)
      \end{itemize}
    \end{column}
    \begin{column}{0.42\textwidth}
      \centering
      \renewcommand{\arraystretch}{1.15}
      \textbf{Proposed $U_w$ / kV}\\[0.3em]
      \begin{tabular}{@{}rr@{}}
        \hline
        550  & 1300 \\
        650  & 1425 \\
        750  & 1550 \\
        850  & 1675 \\
        950  & 1800 \\
        1050 & 1950 \\
        1175 &      \\
        \hline
      \end{tabular}
    \end{column}
  \end{columns}
\end{frame}

\begin{frame}{Example for insulation point coordination}
  Arrester: DC line/cable arrester (DL)
  \vspace{0.5em}
  \begin{columns}[T,onlytextwidth]
    \begin{column}{0.48\textwidth}
      Arrester design:
      \newline\newline
      % \renewcommand{\arraystretch}{1.3}
      % \begin{tabular}{@{}p{4.2cm}r@{}}
      \begin{tabular}{@{}l l@{}}
        Parameter & Value \\
        \hline
        CCOV & $515kV$ \\
        Columns & $8$ \\
        Energy capability & $17MJ$ \\
      \end{tabular}
    \end{column}
    \begin{column}{0.48\textwidth}
      System analysis:
      \newline\newline
      \vspace{1em}
      % \renewcommand{\arraystretch}{1.3}
      % \begin{tabular}{@{}p{4.4cm}r@{}}
      \begin{tabular}{@{}l l@{}}
        Parameter & Value\\
        \hline
        SIPL ($U_{rp}$ for SI) & $807kV @ 1kA$ \\
        LIPL ($U_{rp}$ for LI) & $872kV @ 5kA$ \\
      \end{tabular}
    \end{column}
  \end{columns}
  \vspace{1em}
  Required Switching/Lightning Impulse withstand voltage (RSIWV/ RLIWV):\\
  \begin{tabular}{l l}
    Switching: & $U_{rw} = 1.15 \cdot 807kV = 928kV$\\
    Lightning: & $U_{rw} = 1.2 \cdot 872kV = 1046kV$\\    
  \end{tabular}
  \newline\newline
  Switching/Lightning Impulse withstand voltage (SIWV/ LIWV):\\
    \begin{tabular}{l l}
    Switching: & $U_{w} = 950kV$\\
    Lightning: & $U_{w} = 1050kV$\\    
  \end{tabular}
\end{frame}
